\section{Robustez del estimador y persistencia del tratamiento}
\label{anexo:robustez}

\subsection{Estimador escalonado}

El diseño tiene veinticinco cohortes semestrales de entrada repartidas sobre
doce años. En ese contexto el estimador de diferencias en diferencias con
efectos fijos bidireccionales es un promedio ponderado de comparaciones entre
cohortes en el que unidades ya tratadas actúan como control de unidades tratadas
más tarde, y las ponderaciones pueden ser negativas cuando el efecto del
tratamiento varía entre cohortes \parencite{SUNABRAHAM2021}. Es el problema que
la literatura reciente ha documentado para diseños escalonados.

La \Cref{tab:sunab} reestima el efecto promedio con el estimador ponderado por
interacción de \textcite{SUNABRAHAM2021}, sobre la misma muestra y con los
mismos efectos fijos, de modo que la única diferencia entre columnas es el
estimador.

\begin{table}[H]
\centering
\caption{Efecto promedio de la entrada: Sun y Abraham frente a TWFE}
\label{tab:sunab}
\tabla{tab_sunab}
\end{table}

Las dos columnas son prácticamente idénticas en los cuatro combustibles, y las
diferencias entre ellas son muy inferiores a los errores estándar. La
\Cref{fig:sunab} muestra que la coincidencia se extiende a toda la trayectoria
del estudio de eventos y no solo al promedio.

\begin{figure}[H]
\centering
\includegraphics[width=\textwidth]{sunab_vs_twfe.pdf}
\caption{Trayectoria estimada por Sun y Abraham frente a TWFE}
\label{fig:sunab}
\end{figure}

El resultado no es sorprendente y tiene una explicación conocida: el sesgo por
heterogeneidad de efectos entre cohortes se vuelve pequeño cuando la mayoría de
las unidades de comparación nunca son tratadas, ya que en ese caso las
comparaciones problemáticas —entre una cohorte tratada y otra ya tratada—
reciben poco peso. Este diseño está en esa situación, con cerca de seiscientas
estaciones nunca tratadas frente a algo más de setecientas tratadas.

\subsection{Persistencia del tratamiento}

El diseño trata a la entrada como un aumento discreto y duradero del número de
competidores del mercado local. La \Cref{fig:persistencia} verifica ambos
supuestos con dos estudios de eventos que replican la figura 2 de
\textcite{FISCHER2025}, usando como resultado el propio conteo de estaciones y
no el precio.

\begin{figure}[H]
\centering
\includegraphics[width=\textwidth]{persistencia.pdf}
\caption{Efecto de la entrada sobre el número de estaciones del mercado local}
\label{fig:persistencia}
\end{figure}

El panel izquierdo cuenta todas las estaciones activas dentro de 2 kilómetros
de la incumbente, incluida la entrante. El conteo salta prácticamente en una
unidad en el semestre de la entrada y se mantiene elevado durante los semestres
siguientes: el tratamiento ocurre y perdura. El salto no es exactamente igual a
uno porque no todas las estaciones informan precios cada mes, de modo que el
conteo mide presencia efectiva en el registro y no capacidad instalada.

El panel derecho es el sustantivo. Cuenta solo las estaciones que ya operaban
antes de la entrada, es decir, excluye a la entrante. Esa serie permanece plana:
la entrada no desplaza a las incumbentes del mercado, al menos dentro del
horizonte analizado. La verificación importa porque, de haber desplazamiento, el
efecto estimado sobre los precios sería en parte un efecto de composición
—cambiaría el conjunto de estaciones sobre el que se promedia— y no una
respuesta competitiva de las estaciones que permanecen.

Se observa una leve caída de ambas series en el último bin, que agrupa todo lo
ocurrido más allá de dos años. Es consistente con la atrición natural del
registro a horizontes largos y no altera la lectura de los primeros cuatro
semestres, que es el horizonte sobre el que se interpreta el efecto.
